\documentclass[11pt,a4paper]{article}
\usepackage[utf8]{inputenc}
\usepackage[T1]{fontenc}
\usepackage[french]{babel}
\usepackage[margin=2.5cm]{geometry}
\usepackage{graphicx} % Required for inserting images
\usepackage{float}      % Permet de fixer la figure avec [H]
\usepackage{longtable}
\usepackage{array}
\usepackage{booktabs}
\usepackage[table]{xcolor}
\usepackage{hyperref}
\definecolor{entete}{RGB}{31,78,121}

\title{Projet CloudOps Manager}
\author{Rami Moulehi \and Aziz Ktiti \and Malak Habchi \and Samar Sifaoui}
\date{Octobre 2026}

\newcolumntype{L}[1]{>{\raggedright\arraybackslash}p{#1}}
\newcolumntype{C}[1]{>{\centering\arraybackslash}p{#1}}

% Tableau de besoins fonctionnels réutilisable
% Argument : légende du tableau
\newenvironment{tablebesoins}[1]
{\renewcommand{\arraystretch}{1.3}
 \begin{longtable}{|L{0.09\linewidth}|L{0.48\linewidth}|L{0.16\linewidth}|C{0.12\linewidth}|}
 \hline
\rowcolor{entete}
 \textcolor{white}{\textbf{ID}} & \textcolor{white}{\textbf{Besoin}} & \textcolor{white}{\textbf{Acteur(s)}} & \textcolor{white}{\textbf{Priorité}} \\ \hline
 \endhead
 \caption{#1}
 \endlastfoot}
{\end{longtable}}


\begin{document}

\maketitle
\tableofcontents
\newpage

\section{Spécification des besoins}

\subsection{Introduction}

CloudOps Manager est une plateforme web full-stack qui centralise la gestion des
projets, des serveurs et des applications cloud d'une entreprise. Elle permet aussi
de superviser ces ressources et de traiter les alertes.

Cette partie recense les besoins du système. Elle identifie d'abord les acteurs,
puis présente les besoins fonctionnels par module et les besoins non fonctionnels,
avant de décrire la stack technologique, l'architecture et le diagramme de cas d'utilisation.

Chaque besoin fonctionnel porte un identifiant unique (BF-xx) et une priorité :
\textbf{Haute} pour les fonctionnalités obligatoires du cœur métier,
\textbf{Moyenne} pour celles qui les complètent, \textbf{Basse} pour les
améliorations facultatives (bonus du cahier des charges).

% ---------------------------------------------------------------------
\subsection{Identification des acteurs}

Le système fait intervenir trois acteurs humains authentifiés, qui correspondent aux
trois rôles imposés, un visiteur non authentifié et un acteur système.

\begin{table}[h]
\centering
\renewcommand{\arraystretch}{1.3}
\begin{tabular}{|L{0.22\linewidth}|L{0.14\linewidth}|L{0.50\linewidth}|}
\hline
\rowcolor{entete}
\textcolor{white}{\textbf{Acteur}} & \textcolor{white}{\textbf{Type}} & \textcolor{white}{\textbf{Description}} \\
\hline
Visiteur & Humain & Personne non authentifiée : elle peut uniquement s'inscrire
ou se connecter. \\ \hline
Administrateur & Humain & Gère les utilisateurs, leurs rôles et l'ensemble des
ressources. Il dispose de tous les droits. \\ \hline
DevOps & Humain & Gère les serveurs et les applications, configure les seuils
d'alerte et traite les incidents. \\ \hline
Viewer & Humain & Consulte les projets, les ressources, le tableau de bord et le
monitoring, sans pouvoir rien modifier, à l'exception de son propre profil. \\ \hline
Système (planificateur) & Automatique & Génère ou collecte périodiquement les
mesures, compare les valeurs aux seuils et déclenche les alertes. \\
\hline
\end{tabular}
\caption{Acteurs du système}
\label{tab:acteurs}
\end{table}

% ---------------------------------------------------------------------
\subsection{Besoins fonctionnels}

Les besoins fonctionnels décrivent les services que la plateforme doit rendre à ses
utilisateurs. Ils sont regroupés selon les huit modules obligatoires du cahier des
charges. Dans les tableaux, « Tous » désigne les trois rôles authentifiés
(Admin, DevOps, Viewer).

\subsubsection{Authentification et gestion des utilisateurs}

\begin{tablebesoins}{Besoins fonctionnels -- Authentification et utilisateurs}
BF-01 & Le système doit permettre à un visiteur de s'inscrire, ou à
l'administrateur de créer un compte utilisateur. & Visiteur, Admin & Haute \\ \hline
BF-02 & Le système doit permettre à un utilisateur de s'authentifier par email et
mot de passe et lui délivrer un jeton JWT. & Visiteur & Haute \\ \hline
BF-03 & Le système doit permettre à un utilisateur de se déconnecter (suppression
du jeton côté client). & Tous & Haute \\ \hline
BF-04 & Le système doit permettre à un utilisateur de consulter et modifier son
profil. & Tous & Moyenne \\ \hline
BF-05 & Le système doit permettre à l'administrateur de lister, créer, modifier,
désactiver et supprimer des utilisateurs (sauf son propre compte). & Admin & Haute \\ \hline
BF-06 & Le système doit permettre à l'administrateur d'attribuer un rôle (Admin,
DevOps, Viewer) à chaque utilisateur. & Admin & Haute \\ \hline
BF-07 & Le système doit restreindre l'accès aux routes et aux actions selon le rôle
de l'utilisateur connecté. & Système & Haute \\ \hline
\end{tablebesoins}

\subsubsection{Gestion des projets}

\begin{tablebesoins}{Besoins fonctionnels -- Gestion des projets}
BF-08 & Le système doit permettre de créer, consulter, modifier et supprimer un
projet. & Admin & Haute \\ \hline
BF-09 & Le système doit permettre de créer des équipes et d'y ajouter ou retirer des
membres. & Admin & Haute \\ \hline
BF-10 & Le système doit permettre d'affecter des équipes ou des utilisateurs à un
projet. & Admin & Haute \\ \hline
BF-11 & Le système doit permettre de définir les environnements d'un projet
(développement, staging, production). & Admin, DevOps & Haute \\ \hline
BF-12 & Le système doit afficher la liste des projets avec pagination, recherche et
filtres (statut, équipe, environnement). & Tous & Moyenne \\ \hline
\end{tablebesoins}

\subsubsection{Gestion des serveurs}

\begin{tablebesoins}{Besoins fonctionnels -- Gestion des serveurs}
BF-13 & Le système doit permettre de créer, consulter, modifier et supprimer un
serveur VPS ou VM. & Admin, DevOps & Haute \\ \hline
BF-14 & Le système doit enregistrer pour chaque serveur son fournisseur, sa région,
son CPU, sa RAM, son stockage, son adresse IP et son statut. & Admin, DevOps & Haute \\ \hline
BF-15 & Le système doit permettre de rattacher un serveur à un projet et à un
environnement. & Admin, DevOps & Haute \\ \hline
BF-16 & Le système doit afficher la liste des serveurs avec pagination, recherche et
filtres (fournisseur, région, statut). & Tous & Moyenne \\ \hline
\end{tablebesoins}

\subsubsection{Gestion des applications}

\begin{tablebesoins}{Besoins fonctionnels -- Gestion des applications}
BF-17 & Le système doit permettre de créer, consulter, modifier et supprimer une
application. & Admin, DevOps & Haute \\ \hline
BF-18 & Le système doit permettre d'associer chaque application à un serveur.
& Admin, DevOps & Haute \\ \hline
BF-19 & Le système doit enregistrer pour chaque application sa version, son
environnement et son statut. & Admin, DevOps & Haute \\ \hline
BF-20 & Le système doit afficher la liste des applications avec pagination,
recherche et filtres (serveur, environnement, statut). & Tous & Moyenne \\ \hline
\end{tablebesoins}

\subsubsection{Monitoring}

\begin{tablebesoins}{Besoins fonctionnels -- Monitoring}
BF-21 & Le système doit générer ou collecter périodiquement les mesures CPU et RAM
de chaque serveur. & Système & Haute \\ \hline
BF-22 & Le système doit conserver l'historique horodaté des mesures. & Système & Haute \\ \hline
BF-23 & Le système doit calculer et afficher la disponibilité de chaque serveur.
& Système, Tous & Moyenne \\ \hline
BF-24 & Le système doit afficher des graphiques d'évolution CPU/RAM sur une période
choisie. & Tous & Haute \\ \hline
BF-25 & Le système doit mettre à jour automatiquement l'affichage des mesures.
& Système & Moyenne \\ \hline
\end{tablebesoins}

\subsubsection{Alertes et incidents}

\begin{tablebesoins}{Besoins fonctionnels -- Alertes et incidents}
BF-26 & Le système doit permettre de configurer des seuils d'alerte (par exemple
CPU $>$ 85\,\%) par serveur ou de façon globale. & Admin, DevOps & Haute \\ \hline
BF-27 & Le système doit créer automatiquement une alerte lorsqu'une mesure dépasse
un seuil. & Système & Haute \\ \hline
BF-28 & Le système doit attribuer à chaque alerte un niveau de gravité (info,
avertissement, critique). & Système & Haute \\ \hline
BF-29 & Le système doit permettre de prendre en charge une alerte puis de la marquer
comme résolue. & Admin, DevOps & Moyenne \\ \hline
BF-30 & Le système doit afficher l'historique des incidents, filtrable par gravité,
serveur, statut et période. & Tous & Haute \\ \hline
\end{tablebesoins}

\subsubsection{Tableau de bord}

\begin{tablebesoins}{Besoins fonctionnels -- Tableau de bord}
BF-31 & Le système doit afficher des compteurs globaux : projets, serveurs,
applications et alertes actives. & Tous & Haute \\ \hline
BF-32 & Le système doit afficher des indicateurs clés : serveurs en ligne et hors
ligne, utilisation moyenne CPU/RAM, disponibilité. & Tous & Haute \\ \hline
BF-33 & Le système doit afficher des graphiques de synthèse (répartition par statut
ou fournisseur, évolution des alertes). & Tous & Moyenne \\ \hline
BF-34 & Le système doit permettre de filtrer le tableau de bord par projet,
environnement et période. & Tous & Moyenne \\ \hline
\end{tablebesoins}

\subsubsection{Déploiement}

Ce module ne correspond pas à une interaction d'un utilisateur final : ses besoins
sont des exigences portées par l'équipe de développement.

\begin{tablebesoins}{Besoins fonctionnels -- Déploiement}
BF-35 & L'application complète (frontend, backend, base de données) doit se lancer
avec Docker Compose. & Équipe projet & Haute \\ \hline
BF-36 & La configuration doit passer par des variables d'environnement.
& Équipe projet & Haute \\ \hline
BF-37 & Les données doivent être conservées après redémarrage grâce à un volume
persistant. & Équipe projet & Haute \\ \hline
BF-38 & L'application doit être déployée sur un VPS, une VM ou le cloud et
accessible pour la démonstration. & Équipe projet & Haute \\ \hline
\end{tablebesoins}

\subsubsection{Fonctionnalités bonus}

Le cahier des charges propose des fonctionnalités facultatives, qui peuvent
rapporter jusqu'à deux points de bonus.

\begin{tablebesoins}{Besoins fonctionnels -- Fonctionnalités bonus}
BF-39 & Le système doit pousser les nouvelles mesures et alertes en temps réel
vers l'interface (WebSocket). & Système & Basse \\ \hline
BF-40 & Chaque modification du dépôt doit déclencher automatiquement les tests et
la construction des images (CI/CD avec GitHub Actions). & Équipe projet & Basse \\ \hline
BF-41 & L'application peut être déployée sur un fournisseur cloud (AWS ou Azure).
& Équipe projet & Basse \\ \hline
\end{tablebesoins}

% ---------------------------------------------------------------------

\subsection{Besoins non fonctionnels}

Les besoins non fonctionnels fixent les critères de qualité et les contraintes que
le système doit respecter.

\begin{center}
\renewcommand{\arraystretch}{1.3}
\begin{longtable}{|L{0.10\linewidth}|L{0.18\linewidth}|L{0.58\linewidth}|}
\hline
\rowcolor{entete}
\textcolor{white}{\textbf{ID}} & \textcolor{white}{\textbf{Catégorie}} & \textcolor{white}{\textbf{Besoin}} \\
\hline
\endhead
\caption{Besoins non fonctionnels}
\label{tab:bnf}
\endlastfoot
BNF-01 & Sécurité & Les mots de passe sont stockés hachés (bcrypt) et les routes
protégées exigent un jeton JWT valide. \\ \hline
BNF-02 & Sécurité & Le contrôle d'accès par rôle est appliqué côté backend,
indépendamment de l'interface. \\ \hline
BNF-03 & Fiabilité & Toutes les entrées sont validées par des DTO et des pipes ; les
erreurs sont gérées de façon centralisée avec des réponses cohérentes. \\ \hline
BNF-04 & Performance & Les listes sont paginées pour limiter le volume de données
renvoyé. \\ \hline
BNF-05 & Ergonomie & L'interface est responsive (mobile, tablette, ordinateur) et
affiche les états de chargement et d'erreur. \\ \hline
BNF-06 & Maintenabilité & Le code est organisé en modules NestJS et typé en
TypeScript de bout en bout. \\ \hline
BNF-07 & Interopérabilité & L'API REST expose au moins 20 endpoints documentés avec
Swagger. \\ \hline
BNF-08 & Qualité & Le projet comporte des tests unitaires et quelques tests
d'intégration avec Jest. \\ \hline
BNF-09 & Portabilité & L'application est conteneurisée avec Docker et se déploie de
façon reproductible. \\ \hline
BNF-10 & Persistance & Les données sont stockées dans PostgreSQL et survivent au
redémarrage des conteneurs. \\ \hline
BNF-11 & Traçabilité & L'historique Git reflète la contribution de chaque membre de
l'équipe. \\ \hline
BNF-12 & Conformité & Les mesures de monitoring sont simulées ; une collecte
réelle n'est réalisée que sur une infrastructure de test contrôlée. \\ \hline
BNF-13 & Données & La base de données comporte au moins huit entités liées entre
elles. \\ \hline
\end{longtable}
\end{center}

% ---------------------------------------------------------------------
\subsection{Stack technologique}

Le cahier des charges impose la stack suivante. Le tableau~\ref{tab:stack} précise,
pour chaque couche, les technologies retenues et leur rôle dans CloudOps Manager.

\begin{table}[H]
\centering
\renewcommand{\arraystretch}{1.3}
\begin{tabular}{|L{0.17\linewidth}|L{0.28\linewidth}|L{0.45\linewidth}|}
\hline
\rowcolor{entete}
\textcolor{white}{\textbf{Couche}} & \textcolor{white}{\textbf{Technologies}} & \textcolor{white}{\textbf{Rôle dans le projet}} \\
\hline
Frontend & Next.js, TypeScript, App Router & Interface web responsive : pages
d'authentification, tableau de bord, gestion des ressources et graphiques de
monitoring. \\ \hline
Backend & NestJS, TypeScript, API REST, Swagger & API organisée en modules
(Auth, Projets, Serveurs, Applications, Monitoring, Alertes, Dashboard),
documentée avec Swagger. \\ \hline
Données & PostgreSQL, Prisma & Base relationnelle et persistante. Prisma (ORM)
gère le schéma, les migrations et les requêtes typées\,\footnotemark. \\ \hline
Sécurité & JWT, Guards, contrôle d'accès par rôle & Authentification par jeton,
Guards NestJS et décorateur de rôles (Admin, DevOps, Viewer) sur chaque
route protégée. \\ \hline
Infrastructure & Docker, Docker Compose, VPS / VM / Cloud & Un conteneur par
service (frontend, backend, base de données), orchestrés par Docker Compose et
déployés sur un serveur accessible. \\ \hline
Tests & Jest & Tests unitaires des services et tests d'intégration des
principaux endpoints. \\ \hline
\end{tabular}
\caption{Stack technologique de CloudOps Manager}
\label{tab:stack}
\end{table}
\footnotetext{Le cahier des charges autorise Prisma ou TypeORM. Nous avons choisi
Prisma pour son schéma déclaratif unique, ses migrations automatiques et le typage
TypeScript généré pour les requêtes.}

% ---------------------------------------------------------------------
\newpage
\subsection{Architecture}

L'architecture de CloudOps Manager comprend un frontend Next.js, un backend
NestJS et une base de données PostgreSQL, conteneurisés avec Docker Compose
(figure~\ref{fig:Architecture}).

\begin{figure}[H]
    \centering
    \includegraphics[width=1\linewidth]{Architecture.png}
    \caption{Architecture globale de CloudOps Manager}
    \label{fig:Architecture}
\end{figure}

% ---------------------------------------------------------------------
\subsection{Diagramme de cas d'utilisation}

La figure~\ref{fig:cas-utilisation} présente les cas d'utilisation de la plateforme
et les acteurs qui les déclenchent. Les trois rôles humains sont liés par une relation
de généralisation : l'Administrateur dispose de tous les cas du DevOps, qui dispose
lui-même de tous ceux du Viewer. Le visiteur peut uniquement s'inscrire et
s'authentifier. Le planificateur collecte périodiquement les mesures et génère une
alerte lorsqu'un seuil est dépassé.

\begin{figure}[H]
    \centering
    \includegraphics[width=\linewidth]{diagramme_cas_utilisation.png}
    \caption{Diagramme de cas d'utilisation global de CloudOps Manager}
    \label{fig:cas-utilisation}
\end{figure}

\end{document}
